%% Generated by Sphinx.
\def\sphinxdocclass{report}
\documentclass[a4paper,12pt,spanish]{sphinxmanual}
\ifdefined\pdfpxdimen
   \let\sphinxpxdimen\pdfpxdimen\else\newdimen\sphinxpxdimen
\fi \sphinxpxdimen=.75bp\relax


\catcode`^^^^00a0\active\protected\def^^^^00a0{\leavevmode\nobreak\ }
\usepackage{cmap}
\usepackage{fontspec}
\usepackage{amsmath,amssymb,amstext}
\usepackage{polyglossia}
\setmainlanguage{spanish}

        \usepackage{libertine}
        \setmonofont[
            Path=../../fonts/,
            BoldFont=FiraMono-Bold.ttf,
            AutoFakeSlant,
            BoldItalicFeatures={FakeSlant},
            Scale=MatchLowercase
        ]{FiraMono-Medium.ttf}
\usepackage[Sonny]{fncychap}
\usepackage[dontkeepoldnames]{sphinx}
\sphinxsetup{
        verbatimwithframe=false,
        noteBorderColor={rgb}{0.208,0.374,0.486},
        noteborder=1pt}
\usepackage{geometry}

% Include hyperref last.
\usepackage{hyperref}
% Fix anchor placement for figures with captions.
\usepackage{hypcap}% it must be loaded after hyperref.
% Set up styles of URL: it should be placed after hyperref.
\urlstyle{same}

\addto\captionsspanish{\renewcommand{\figurename}{Figura}}
\addto\captionsspanish{\renewcommand{\tablename}{Tabla}}
\addto\captionsspanish{\renewcommand{\literalblockname}{Lista}}

\addto\captionsspanish{\renewcommand{\literalblockcontinuedname}{continued from previous page}}
\addto\captionsspanish{\renewcommand{\literalblockcontinuesname}{continues on next page}}

\def\pageautorefname{página}

\setcounter{tocdepth}{0}

\usepackage[dotinlabels]{titletoc}
\usepackage{titlesec}
\usepackage{array,tabularx}
\usepackage{colortbl}
\usepackage{xcolor,graphicx}
\usepackage[xparse,skins,breakable]{tcolorbox}
\titlelabel{\thetitle.\quad}
\let\sphinxcodeORI\sphinxcode
\protected\def\sphinxcode #1%
    {{\color{OuterLinkColor}\sphinxcodeORI{#1}}}

\DeclareTColorBox{caja}{ m m O{} }
    { center,width=1\linewidth,enhanced,fonttitle=\large\bfseries,
      drop fuzzy shadow,breakable,adjusted title={#1},
      before skip=0.5cm,after skip=0.5cm,
      bottom=0.4cm,
      left=0.5cm,
      right=0.5cm,
      before upper={ \vskip -1.8cm\hfill%
                     \includegraphics[width=0.13\textwidth]{../../iconos/#2}%
                     \vskip -0.1cm\vfill%
                     \small%
                   },
      #3
    }

\renewenvironment{sphinxnote}[1]
    {\begin{caja}{#1}{bulb.png}}
    {\end{caja}}

\renewenvironment{sphinxdanger}[1]
    {\begin{caja}{#1}{emblem-ohno.png}[colframe=red!75!black]}
    {\end{caja}}

%\renewenvironment{sphinxnote}[1]
%    {\begin{tcolorbox}[center,width=1\linewidth,enhanced,fonttitle=\large\bfseries,drop
%    fuzzy shadow,breakable,adjusted title=#1,before skip=0.5cm,after skip=0.5cm]
%     \vskip -2cm\hfill%
%     \includegraphics[width=0.15\textwidth]{../../iconos/hint.png}%
%     \vskip -0.1cm\vfill%
%     \small
%    }
%    {\end{tcolorbox}}


\title{Desarrollo web en entorno servidor}
\date{enero de 2018}
\release{1.0}
\author{Ricardo Pérez López}
\newcommand{\sphinxlogo}{\vbox{}}
\renewcommand{\releasename}{Versión}
\makeindex

\begin{document}

\maketitle
\sphinxtableofcontents
\phantomsection\label{\detokenize{index_latex::doc}}



\part{Introducción}
\label{\detokenize{introduccion:introduccion}}\label{\detokenize{introduccion:desarrollo-web-en-entorno-servidor}}\label{\detokenize{introduccion::doc}}

\chapter{Preparación del entorno de desarrollo}
\label{\detokenize{introduccion:preparacion-del-entorno-de-desarrollo}}

\section{Instalación automatizada}
\label{\detokenize{introduccion:instalacion-automatizada}}

\subsection{Acciones previas}
\label{\detokenize{introduccion:acciones-previas}}

\subsubsection{Instalar \sphinxstyleliteralintitle{git}}
\label{\detokenize{introduccion:instalar-git}}

\subsubsection{Crear cuenta en GitHub}
\label{\detokenize{introduccion:crear-cuenta-en-github}}

\subsubsection{Usar https://github.com/ricpelo/conf y seguir las instrucciones del README.md}
\label{\detokenize{introduccion:usar-https-github-com-ricpelo-conf-y-seguir-las-instrucciones-del-readme-md}}

\section{Terminal}
\label{\detokenize{introduccion:terminal}}

\subsection{Zsh}
\label{\detokenize{introduccion:zsh}}

\subsection{Oh My Zsh}
\label{\detokenize{introduccion:oh-my-zsh}}

\section{Navegador}
\label{\detokenize{introduccion:navegador}}

\subsection{Herramientas de desarrollador}
\label{\detokenize{introduccion:herramientas-de-desarrollador}}

\section{Editores de texto}
\label{\detokenize{introduccion:editores-de-texto}}

\subsection{Vim y less}
\label{\detokenize{introduccion:vim-y-less}}

\subsection{Atom}
\label{\detokenize{introduccion:atom}}

\chapter{Introducción al desarrollo web}
\label{\detokenize{introduccion:introduccion-al-desarrollo-web}}

\section{Conceptos básicos}
\label{\detokenize{introduccion:conceptos-basicos}}

\subsection{Navegadores y servidores web}
\label{\detokenize{introduccion:navegadores-y-servidores-web}}

\subsection{Agentes de usuario}
\label{\detokenize{introduccion:agentes-de-usuario}}

\subsection{Web estática vs. dinámica}
\label{\detokenize{introduccion:web-estatica-vs-dinamica}}

\subsection{Estructura vs. contenido}
\label{\detokenize{introduccion:estructura-vs-contenido}}

\section{Ejemplos de aplicaciones web}
\label{\detokenize{introduccion:ejemplos-de-aplicaciones-web}}

\subsection{Redes sociales: Facebook, Twitter…}
\label{\detokenize{introduccion:redes-sociales-facebook-twitter}}

\subsection{Comercio electrónico: Amazon, eBay…}
\label{\detokenize{introduccion:comercio-electronico-amazon-ebay}}

\subsection{Administración electrónica…}
\label{\detokenize{introduccion:administracion-electronica}}

\subsection{Portales}
\label{\detokenize{introduccion:portales}}

\subsection{ERP, CRM}
\label{\detokenize{introduccion:erp-crm}}

\section{Tecnologías de desarrollo de aplicaciones web}
\label{\detokenize{introduccion:tecnologias-de-desarrollo-de-aplicaciones-web}}

\subsection{.NET}
\label{\detokenize{introduccion:net}}

\subsection{Java}
\label{\detokenize{introduccion:java}}

\subsection{Ruby/Rails}
\label{\detokenize{introduccion:ruby-rails}}

\subsection{Python/Django}
\label{\detokenize{introduccion:python-django}}

\subsection{PHP}
\label{\detokenize{introduccion:php}}

\subsection{El Kung-Fu de la programación}
\label{\detokenize{introduccion:el-kung-fu-de-la-programacion}}

\subsubsection{Odoo}
\label{\detokenize{introduccion:odoo}}

\subsubsection{PrestaShop}
\label{\detokenize{introduccion:prestashop}}

\subsubsection{Drupal}
\label{\detokenize{introduccion:drupal}}

\subsubsection{WordPress}
\label{\detokenize{introduccion:wordpress}}

\chapter{Protocolo HTTP y lenguaje HTML}
\label{\detokenize{introduccion:protocolo-http-y-lenguaje-html}}

\section{Arquitectura cliente/servidor}
\label{\detokenize{introduccion:arquitectura-cliente-servidor}}

\section{HTML 5 básico}
\label{\detokenize{introduccion:html-5-basico}}

\section{Protocolo HTTP}
\label{\detokenize{introduccion:protocolo-http}}

\subsection{URIs}
\label{\detokenize{introduccion:uris}}

\subsubsection{URL encoding}
\label{\detokenize{introduccion:url-encoding}}

\subsection{Peticiones y respuestas}
\label{\detokenize{introduccion:peticiones-y-respuestas}}

\subsection{Métodos: GET y POST}
\label{\detokenize{introduccion:metodos-get-y-post}}

\subsection{Cabeceras HTTP}
\label{\detokenize{introduccion:cabeceras-http}}

\subsection{Códigos de estado}
\label{\detokenize{introduccion:codigos-de-estado}}

\subsection{Experimentos}
\label{\detokenize{introduccion:experimentos}}

\subsubsection{\sphinxstyleliteralintitle{telnet}}
\label{\detokenize{introduccion:telnet}}

\subsubsection{\sphinxstyleliteralintitle{netcat}}
\label{\detokenize{introduccion:netcat}}

\subsubsection{\sphinxstyleliteralintitle{curl}}
\label{\detokenize{introduccion:curl}}

\subsubsection{\sphinxstyleliteralintitle{http}}
\label{\detokenize{introduccion:http}}

\subsubsection{Google Chrome Developer Tools}
\label{\detokenize{introduccion:google-chrome-developer-tools}}

\subsection{Envío de datos al servidor}
\label{\detokenize{introduccion:envio-de-datos-al-servidor}}

\subsubsection{Mediante GET}
\label{\detokenize{introduccion:mediante-get}}

\subsubsection{Mediante POST}
\label{\detokenize{introduccion:mediante-post}}

\subsubsection{Formularios HTML}
\label{\detokenize{introduccion:formularios-html}}

\subsection{Cookies}
\label{\detokenize{introduccion:cookies}}

\section{Apache básico}
\label{\detokenize{introduccion:apache-basico}}

\subsection{Instalación}
\label{\detokenize{introduccion:instalacion}}

\subsection{Configuración básica}
\label{\detokenize{introduccion:configuracion-basica}}

\subsection{Sitios virtuales}
\label{\detokenize{introduccion:sitios-virtuales}}

\section{Scripts CGI}
\label{\detokenize{introduccion:scripts-cgi}}

\subsection{Configuración de Apache}
\label{\detokenize{introduccion:configuracion-de-apache}}

\subsection{Ejemplos en Ruby}
\label{\detokenize{introduccion:ejemplos-en-ruby}}

\chapter{Sistemas de control de versiones}
\label{\detokenize{introduccion:sistemas-de-control-de-versiones}}

\section{Git básico}
\label{\detokenize{introduccion:git-basico}}\begin{enumerate}
\item {} 
config, init, add, commit

\item {} 
status, log, diff

\item {} 
Alias lg

\item {} 
checkout, reset, revert, \textendash{}amend

\item {} 
show

\item {} 
rm, mv

\item {} 
Atom y Git

\end{enumerate}


\section{Git avanzado}
\label{\detokenize{introduccion:git-avanzado}}\begin{enumerate}
\item {} 
Ramas: branch

\item {} 
merge, rebase

\item {} 
Resolución de conflictos

\end{enumerate}


\section{Ramas remotas}
\label{\detokenize{introduccion:ramas-remotas}}\begin{enumerate}
\item {} 
clone, fetch, push, pull

\item {} 
Ramas de seguimiento (tracking branch)

\end{enumerate}


\section{Repositorios de código (Github.com, GitLab.com, Bitbucket.org…)}
\label{\detokenize{introduccion:repositorios-de-codigo-github-com-gitlab-com-bitbucket-org}}

\part{PHP}
\label{\detokenize{php:php}}\label{\detokenize{php::doc}}

\chapter{Conceptos básicos de PHP I}
\label{\detokenize{php:conceptos-basicos-de-php-i}}
ricpelo’s note: Programada inicialmente para empezar el 23-10-2017. \sphinxcode{\_\_\_\_\_\_}

\begin{sphinxVerbatim}[commandchars=\\\{\}]
Hola
\end{sphinxVerbatim}

\begin{sphinxVerbatim}[commandchars=\\\{\}]
\PYGZdl{} php ../../p.php
Hola
\end{sphinxVerbatim}


\section{Introducción a PHP}
\label{\detokenize{php:introduccion-a-php}}
\begin{sphinxadmonition}{danger}{Peligro:}
¡Esto es un peligro!
\end{sphinxadmonition}


\subsection{Página web de PHP}
\label{\detokenize{php:pagina-web-de-php}}
\textless{}\sphinxurl{http://php.net}\textgreater{}


\subsection{Instalación de PHP}
\label{\detokenize{php:instalacion-de-php}}

\subsection{Documentación y búsqueda de información}
\label{\detokenize{php:documentacion-y-busqueda-de-informacion}}

\subsection{Configuración básica con \sphinxstyleliteralintitle{php.ini}}
\label{\detokenize{php:configuracion-basica-con-php-ini}}

\subsubsection{\sphinxstyleliteralintitle{error\_reporting = E\_ALL}}
\label{\detokenize{php:error-reporting-e-all}}

\subsubsection{\sphinxstyleliteralintitle{display\_errors = On}}
\label{\detokenize{php:display-errors-on}}

\subsubsection{\sphinxstyleliteralintitle{display\_startup\_errors = On}}
\label{\detokenize{php:display-startup-errors-on}}

\subsubsection{\sphinxstyleliteralintitle{date.timezone = 'UTC'}}
\label{\detokenize{php:date-timezone-utc}}

\section{Sintaxis básica}
\label{\detokenize{php:sintaxis-basica}}
\textless{}\sphinxurl{http://php.net/manual/es/language.basic-syntax.php}\textgreater{}


\subsection{Datos e instrucciones}
\label{\detokenize{php:datos-e-instrucciones}}
\index{datos}\index{instrucciones}\ignorespaces 
En todo lenguaje de programación existen dos elementos básicos: los \sphinxstylestrong{datos} y
las \sphinxstylestrong{instrucciones}. Los datos son las porciones de información con las que
trabajan los programas, y las instrucciones son las acciones que manipulan esos
datos para llevar a cabo la tarea para la que fue concebido el programa. Ambos
elementos, datos e instrucciones, constituyen los pilares del lenguaje y de los
programas que se escriben con él.

\index{sentencias}\index{expresiones}\ignorespaces 
Desde un punto de vista sintáctico, en el código fuente del programa, los datos
se codifican en forma de \sphinxstylestrong{expresiones}, y las instrucciones toman la forma
de \sphinxstylestrong{sentencias}.

La diferencia fundamental entre una expresión y una sentencia es la siguiente:

\index{expresiones, evaluación de}\index{evaluación|see{expresiones}}\ignorespaces \begin{itemize}
\item {} 
Una expresión tiene un valor (se dice que \sphinxstyleemphasis{denota} o \sphinxstyleemphasis{representa} un valor),
y por eso decimos que una expresión \sphinxstyleemphasis{se puede evaluar}, y al evaluarse, se
obtiene el valor de la expresión, que no es más que un dato.

\item {} 
En cambio, una sentencia no denota ningún valor, por lo que no puede
evaluarse. El intérprete \sphinxstyleemphasis{ejecuta} la sentencia y se llevan a cabo las
acciones que provoca dicha ejecución.

\end{itemize}

Las expresiones se evalúan. Las sentencias se ejecutan.

\begin{sphinxadmonition}{note}{Ejemplos de expresiones}

\begin{sphinxVerbatim}[commandchars=\\\{\}]
\PYG{l+m+mi}{5} \PYG{o}{*} \PYG{p}{(}\PYG{l+m+mi}{3} \PYG{o}{+} \PYG{l+m+mi}{6}\PYG{p}{)}
\end{sphinxVerbatim}

Es una expresión aritmética que involucra números enteros y cuyo valor es
\sphinxcode{45}.

\begin{sphinxVerbatim}[commandchars=\\\{\}]
\PYG{l+m+mf}{1.6} \PYG{o}{+} \PYG{l+m+mf}{2.3}
\end{sphinxVerbatim}

Es una expresión aritmética que involucra números reales y cuyo valor es
\sphinxcode{3.9}.

\begin{sphinxVerbatim}[commandchars=\\\{\}]
\PYG{l+m+mi}{74}
\end{sphinxVerbatim}

Es una constante literal numérica cuyo valor es, precisamente, \sphinxcode{74}.
\end{sphinxadmonition}


\subsection{Sentencias y comandos}
\label{\detokenize{php:sentencias-y-comandos}}
\index{sentencias}\index{comandos}\ignorespaces 
Las sentencias en PHP pueden ser \sphinxstyleemphasis{simples} o \sphinxstyleemphasis{compuestas}.
\begin{itemize}
\item {} 
Las \sphinxstylestrong{sentencias simples} son las instrucciones más elementales del lenguaje
y se escriben siempre acabadas en punto y coma (\sphinxcode{;}).

\item {} 
Las \sphinxstylestrong{sentencias compuestas} corresponden a las \sphinxstylestrong{estructuras de control} y
se estudiarán posteriormente en este capítulo.

\end{itemize}

Se puede construir una sentencia simple usando sencillamente una expresión y
acabándola en punto y coma, como por ejemplo:

\begin{sphinxVerbatim}[commandchars=\\\{\}]
\PYG{l+m+mi}{8} \PYG{o}{+} \PYG{l+m+mi}{3}\PYG{p}{;}
\end{sphinxVerbatim}

Pero una sentencia así no tendría mucha utilidad, ya que el intérprete de PHP
se limitaría a evaluar la expresión pero no haría nada más con el valor
calculado.

\index{efectos laterales}\ignorespaces 
Las sentencias realmente útiles son aquellas que provocan \sphinxstylestrong{efectos
laterales}, es decir, acciones que provocan cambios en el estado interno del
programa o que producen resultados que se vuelcan hacia la \sphinxstyleemphasis{salida} (siendo
esta cualquier dispositivo de salida, como por ejemplo la pantalla, un archivo
del disco o una fila de una tabla de una base de datos relacional).

\index{comandos}\index{palabras clave}\ignorespaces 
Otra forma de construir una sentencia simple es usar \sphinxstylestrong{comandos}. PHP dispone
de varios comandos con los que se pueden escribir sentencias para llevar a cabo
instrucciones sencillas. Cada comando consta de una \sphinxstylestrong{palabra clave}, que
identifica al comando, y de una serie de \sphinxstyleemphasis{argumentos} que completan la
sentencia.


\subsubsection{El comando \sphinxstyleliteralintitle{echo}}
\label{\detokenize{php:el-comando-echo}}
\index{echo}\ignorespaces 
El ejemplo clásico de comando en PHP es \sphinxhref{http://php.net/manual/es/function.echo.php}{\sphinxcode{echo}}%
\begin{footnote}[1]\sphinxAtStartFootnote
\sphinxnolinkurl{http://php.net/manual/es/function.echo.php}
%
\end{footnote}. El comando \sphinxcode{echo} vuelca a
la salida el valor de las expresiones que se indican como parámetro en la
sentencia. Por ejemplo:

\begin{sphinxVerbatim}[commandchars=\\\{\}]
\PYG{k}{echo} \PYG{l+m+mi}{25} \PYG{o}{*} \PYG{l+m+mi}{3}\PYG{p}{;}
\end{sphinxVerbatim}

Muestra \sphinxcode{75} por la salida (normalmente la pantalla). O bien:

\begin{sphinxVerbatim}[commandchars=\\\{\}]
\PYG{k}{echo} \PYG{l+s+s1}{\PYGZsq{}¡Hola a todos!\PYGZsq{}}\PYG{p}{;}
\end{sphinxVerbatim}

Muestra \sphinxcode{¡Hola a todos!}.

Puede mostrar varios valores, separando cada uno de ellos entre sí con una
coma:

\begin{sphinxVerbatim}[commandchars=\\\{\}]
\PYG{k}{echo} \PYG{l+s+s1}{\PYGZsq{}El resultado es: \PYGZsq{}}\PYG{p}{,} \PYG{l+m+mi}{4} \PYG{o}{*} \PYG{l+m+mi}{2}\PYG{p}{;}
\end{sphinxVerbatim}

Mostraría \sphinxcode{El resultado es: 8}.

\index{expresiones}\ignorespaces 

\subsection{Expresiones}
\label{\detokenize{php:expresiones}}\label{\detokenize{php:index-7}}
El otro tipo de construcción sintáctica que existe en PHP además de las
sentencias son las \sphinxstyleemphasis{expresiones}. Una expresión \sphinxstyleemphasis{denota} o \sphinxstyleemphasis{representa} un
valor. Una expresión puede ser tan simple como una constante literal (por
ejemplo, el número \sphinxcode{25}) o tan compleja que involucre constantes, variables,
operadores, funciones, métodos… combinados todos ellos entre sí para formar
una única expresión.


\subsubsection{Operadores}
\label{\detokenize{php:operadores}}
\index{operadores}\ignorespaces 
Un \sphinxstylestrong{operador} es un símbolo que representa una operación que se desea
realizar sobre sus \sphinxstylestrong{operandos} %
\begin{footnote}[2]\sphinxAtStartFootnote
El número de operandos de un operador se denomina \sphinxstylestrong{aridad}. La aridad
puede ser 1, 2 ó 3, según el operador sea \sphinxstyleemphasis{unario}, \sphinxstyleemphasis{binario} o \sphinxstyleemphasis{ternario},
respectivamente.
%
\end{footnote}. Los operandos son los valores
sobre los que actúa el operador para llevar a cabo la operación deseada. Por
ejemplo:

\begin{sphinxVerbatim}[commandchars=\\\{\}]
\PYG{l+m+mi}{4} \PYG{o}{+} \PYG{l+m+mi}{3}
\end{sphinxVerbatim}

Aquí, el operador \sphinxcode{+} representa la operación \sphinxstyleemphasis{suma} a realizar sobre los
números \sphinxcode{4} y \sphinxcode{3}, que son sus operandos. Como el operador actúa sobre dos
operandos, se dice que es un operador \sphinxstyleemphasis{binario}. En cambio:

\begin{sphinxVerbatim}[commandchars=\\\{\}]
\PYG{o}{\PYGZhy{}}\PYG{l+m+mi}{17}
\end{sphinxVerbatim}

Aquí se usa el operador \sphinxcode{-} (\sphinxstyleemphasis{signo menos}) para convertir en negativo el
valor \sphinxcode{17}. Como el operador actúa sobre un único operando, se dice que es un
operador \sphinxstyleemphasis{unario}.

En PHP existe un único operador \sphinxstyleemphasis{ternario} que se estudiará posteriormente.

En una misma expresión pueden actuar varios operadores, como en:

\begin{sphinxVerbatim}[commandchars=\\\{\}]
\PYG{l+m+mi}{4} \PYG{o}{+} \PYG{l+m+mi}{3} \PYG{o}{+} \PYG{l+m+mi}{5}
\end{sphinxVerbatim}

Que denota el valor \sphinxcode{12}, o con varios operadores diferentes:

\begin{sphinxVerbatim}[commandchars=\\\{\}]
\PYG{l+m+mi}{4} \PYG{o}{+} \PYG{l+m+mi}{3} \PYG{o}{*} \PYG{l+m+mi}{5}
\end{sphinxVerbatim}

Que evalúa a \sphinxcode{19}.


\paragraph{Asociatividad y prioridad}
\label{\detokenize{php:asociatividad-y-prioridad}}
\index{asociatividad}\index{prioridad}\ignorespaces 
Todas las expresiones anteriores son ejemplos de expresiones \sphinxstyleemphasis{artiméticas},
donde se realizan las operaciones matemáticas usuales (suma, resta, producto y
división) sobre números. La evaluación de una expresión (ya sea aritmética o de
cualquier otro tipo) depende de las reglas de \sphinxstylestrong{asociatividad} y \sphinxstylestrong{prioridad}
de los operadores que participan en dicha expresión, las cuales tenemos que
conocer para entender cómo evaluará el intérprete las expresiones que formen
parte de nuestro programa. En el caso de las expresiones aritméticas, las
reglas son las habituales que aprendimos en el colegio:
\begin{itemize}
\item {} 
En una expresión en la que un operando está rodeado a izquierda y derecha por
\sphinxstyleemphasis{el mismo operador}, se aplica la regla de la \sphinxstyleemphasis{asociatividad}. Por ejemplo,
en la expresión:

\begin{sphinxVerbatim}[commandchars=\\\{\}]
\PYG{l+m+mi}{4} \PYG{o}{+} \PYG{l+m+mi}{3} \PYG{o}{+} \PYG{l+m+mi}{5}
\end{sphinxVerbatim}

el operando \sphinxcode{3} tiene el mismo operador a izquierda y derecha (el \sphinxcode{+}), y
como dicho operador es \sphinxstyleemphasis{asociativo por la izquierda}, la expresión se evalúa
igual que si se hubiera escrito como:

\begin{sphinxVerbatim}[commandchars=\\\{\}]
\PYG{p}{(}\PYG{l+m+mi}{4} \PYG{o}{+} \PYG{l+m+mi}{3}\PYG{p}{)} \PYG{o}{+} \PYG{l+m+mi}{5}
\end{sphinxVerbatim}

\item {} 
En una expresión en la que un operando está rodeado a izquierda y derecha por
\sphinxstyleemphasis{distintos operadores}, se aplica la regla de la \sphinxstyleemphasis{prioridad}. Por ejemplo,
en la expresión:

\begin{sphinxVerbatim}[commandchars=\\\{\}]
\PYG{l+m+mi}{4} \PYG{o}{+} \PYG{l+m+mi}{3} \PYG{o}{*} \PYG{l+m+mi}{5}
\end{sphinxVerbatim}

el operando \sphinxcode{3} tiene el operador \sphinxcode{+} a su izquierda y el \sphinxcode{*} a su
derecha, pero como el producto tiene más prioridad que la suma, la expresión
se evalúa igual que si se hubiera escrito como:

\begin{sphinxVerbatim}[commandchars=\\\{\}]
\PYG{l+m+mi}{4} \PYG{o}{+} \PYG{p}{(}\PYG{l+m+mi}{3} \PYG{o}{*} \PYG{l+m+mi}{5}\PYG{p}{)}
\end{sphinxVerbatim}

\end{itemize}

Como se aprecia en los ejemplos anteriores, se pueden usar \sphinxstylestrong{paréntesis} para
agrupar sub-expresiones dentro de una expresión y así aumentar la prioridad de
los operadores que vayan entre paréntesis. Por ejemplo, en la expresión:

\begin{sphinxVerbatim}[commandchars=\\\{\}]
\PYG{p}{(}\PYG{l+m+mi}{4} \PYG{o}{+} \PYG{l+m+mi}{3}\PYG{p}{)} \PYG{o}{*} \PYG{l+m+mi}{5}
\end{sphinxVerbatim}

la suma se hace antes que el producto, aunque este último sea un operador de
mayor prioridad. El resultado de dicha expresión es el valor \sphinxcode{35}.


\subsubsection{Funciones}
\label{\detokenize{php:funciones}}
\index{funciones}\index{cos()}\ignorespaces 
Las funciones en las expresiones cumplen el mismo papel que en Matemáticas:
realizan un cálculo a partir de unos valores de entrada indicados en sus
argumentos y \sphinxstyleemphasis{devuelven} el resultado de dicho cálculo. Por ejemplo, la
función \sphinxstyleemphasis{coseno} (abreviada como \sphinxhref{http://php.net/manual/es/function.cos.php}{\sphinxcode{cos()}}%
\begin{footnote}[3]\sphinxAtStartFootnote
\sphinxnolinkurl{http://php.net/manual/es/function.cos.php}
%
\end{footnote}) calcula el coseno de un ángulo. En
Matemáticas (y en Programación) se representa indicando el nombre de la
función y, a continuación, la lista de sus argumentos entre paréntesis y
separados por comas. Así, para calcular el coseno de 2.4 radianes, podemos
escribir:

\begin{sphinxVerbatim}[commandchars=\\\{\}]
\PYG{n+nb}{cos}\PYG{p}{(}\PYG{l+m+mf}{2.4}\PYG{p}{)}
\end{sphinxVerbatim}

Que da como resultado \sphinxcode{-0.73739371554125}, y ese sería el valor de dicha
expresión.

El coseno es un ejemplo de función con un único argumento, pero hay funciones
que admiten o requieren más argumentos. Es el caso de la función \sphinxhref{http://php.net/manual/es/function.max.php}{\sphinxcode{max()}}%
\begin{footnote}[4]\sphinxAtStartFootnote
\sphinxnolinkurl{http://php.net/manual/es/function.max.php}
%
\end{footnote}, que
devuelve el valor máximo de todos los indicados en su lista de argumentos:

\begin{sphinxVerbatim}[commandchars=\\\{\}]
\PYG{n+nb}{max}\PYG{p}{(}\PYG{l+m+mi}{5}\PYG{p}{,} \PYG{l+m+mi}{3}\PYG{p}{,} \PYG{l+m+mi}{8}\PYG{p}{,} \PYG{l+m+mi}{2}\PYG{p}{)}
\end{sphinxVerbatim}

Devuelve \sphinxcode{8}.

\begin{sphinxadmonition}{note}{Nota:}
Cuando usamos una función en una expresión, decimos que estamos \sphinxstyleemphasis{llamando} o
\sphinxstyleemphasis{invocando} a la función. La aparición de la función en la expresión es una
\sphinxstyleemphasis{llamada} a la función.
\end{sphinxadmonition}

\index{efectos laterales}\index{var\_dump()}\ignorespaces 
En PHP, a diferencia de lo que ocurre en Matemáticas, existen funciones que no
devuelven ningún valor, ya que su objetivo es provocar un \sphinxstyleemphasis{efecto lateral}. La
más conocida podría ser, sin temor a equivocarnos, la función \sphinxhref{http://php.net/manual/es/function.var-dump.php}{\sphinxcode{var\_dump()}}%
\begin{footnote}[5]\sphinxAtStartFootnote
\sphinxnolinkurl{http://php.net/manual/es/function.var-dump.php}
%
\end{footnote}. Esta
función muestra en la salida información estructurada sobre las expresiones que
se le pasan como argumento, incluyendo su valor y su tipo. En cierto sentido,
podría considerarse un versión especializada del comando \sphinxcode{echo}, pero en forma
de función y más orientada a la \sphinxstyleemphasis{depuración} de programas.

Es importante destacar que esa información que muestra se vuelca \sphinxstyleemphasis{en la salida}
(normalmente la pantalla). No estamos diciendo que la función \sphinxstyleemphasis{devuelva} dicha
información o que esa información sea el valor resultante de \sphinxstyleemphasis{evaluar} la
llamada a la función. De hecho, estamos hablamos de llamar a la función como si
fuera una sentencia (una sentencia formada únicamente por la llamada a la
función y el punto y coma final):

\begin{sphinxVerbatim}[commandchars=\\\{\}]
\PYG{n+nb}{var\PYGZus{}dump}\PYG{p}{(}\PYG{l+m+mi}{14} \PYG{o}{+} \PYG{l+m+mi}{3}\PYG{p}{);}
\end{sphinxVerbatim}

La sentencia anterior (sí: \sphinxstyleemphasis{sentencia}, porque es una instrucción en sí misma,
terminada en punto y coma), como cualquier otra sentencia, no devuelve ningún
valor, sino que produce un efecto lateral. En este caso, mostrar a la salida
(la pantalla) lo siguiente:

\begin{sphinxVerbatim}[commandchars=\\\{\}]
\PYG{n+nx}{int}\PYG{p}{(}\PYG{l+m+mi}{17}\PYG{p}{)}
\end{sphinxVerbatim}

\index{enteros}\index{números}\ignorespaces 
Observamos que \sphinxcode{var\_dump()} nos informa del valor de la expresión (\sphinxcode{17}) y de
su tipo (\sphinxcode{int}, que significa \sphinxstylestrong{número entero}). La importancia de conocer
el valor y el tipo de toda expresión que aparezca en nuestros programas se
apreciará en breve. Por ahora, veamos dos ejemplos más de utilización de
\sphinxcode{var\_dump()}:

\begin{sphinxVerbatim}[commandchars=\\\{\}]
\PYG{n+nb}{var\PYGZus{}dump}\PYG{p}{(}\PYG{l+m+mf}{12.3} \PYG{o}{\PYGZhy{}} \PYG{l+m+mi}{4}\PYG{p}{);}
\end{sphinxVerbatim}

Muestra a la salida:

\begin{sphinxVerbatim}[commandchars=\\\{\}]
\PYG{n+nx}{float}\PYG{p}{(}\PYG{l+m+mf}{8.3}\PYG{p}{)}
\end{sphinxVerbatim}

\index{reales}\index{números}\ignorespaces 
Lo que nos indica que el valor de la expresión \sphinxcode{12.3 - 4} es el \sphinxstylestrong{número
real} \sphinxcode{8.3}. Finalmente:

\begin{sphinxVerbatim}[commandchars=\\\{\}]
\PYG{n+nb}{var\PYGZus{}dump}\PYG{p}{(}\PYG{l+s+s2}{\PYGZdq{}}\PYG{l+s+s2}{Saludos}\PYG{l+s+s2}{\PYGZdq{}}\PYG{p}{);}
\end{sphinxVerbatim}

Muestra:

\begin{sphinxVerbatim}[commandchars=\\\{\}]
\PYG{n+nx}{string}\PYG{p}{(}\PYG{l+m+mi}{7}\PYG{p}{)} \PYG{l+s+s2}{\PYGZdq{}}\PYG{l+s+s2}{Saludos}\PYG{l+s+s2}{\PYGZdq{}}
\end{sphinxVerbatim}

Que nos informa de que \sphinxcode{"Saludos"} es una cadena de siete caracteres.

\index{tipos de datos}\ignorespaces 
Los números (enteros y reales), así como las cadenas, son algunos de los
\sphinxstylestrong{tipos de datos} que el lenguaje PHP nos proporciona para ayudarnos a
manipular la información. Su estudio detallado se hará posteriormente.


\section{Funcionamiento del intérprete}
\label{\detokenize{php:funcionamiento-del-interprete}}
PHP es un lenguaje \sphinxstyleemphasis{interpretado} y, como tal, requiere de la existencia de un
\sphinxstylestrong{intérprete}, que es la utilidad encargada de leer el código fuente escrito
en el lenguaje y ejecutarlo adecuadamente siguiendo las reglas de dicho
lenguaje.

La ejecución de nuestro código se puede llevar a cabo de dos formas:

\index{scripts}\ignorespaces \begin{description}
\item[{Por lotes:}] \leavevmode
Nuestro código fuente está almacenado en un archivo de texto (normalmente,
con extensión \sphinxcode{.php}) y el intérprete lee dicho archivo, lo analiza
sintáctica y semánticamente y ejecuta las instrucciones que lo forman. Estos
archivos (que contienen el código fuente en PHP) se denominan \sphinxstylestrong{scripts}, y
el objetivo final de este curso es desarrollar aplicaciones escribiendo los
\sphinxstyleemphasis{scripts} necesarios para cumplir con la funcionalidad deseada.

Una manera de ejecutar nuestro \sphinxstyleemphasis{script} es pasárselo al intérprete desde la
consola del sistema operativo %
\begin{footnote}[6]\sphinxAtStartFootnote
La otra forma es ejecutar el \sphinxstyleemphasis{script} en el contexto de un \sphinxstylestrong{servidor
web}, \sphinxstyleemphasis{embebiendo} (o \sphinxstyleemphasis{incrustando}) el código PHP en una página HTML.
Hablaremos sobre ello en próximos capítulos.
%
\end{footnote}. Por ejemplo, si tenemos nuestro
\sphinxstyleemphasis{script} almacenado en el archivo \sphinxcode{prueba.php}, podemos provocar la
ejecución del mismo mediante la siguiente orden del sistema operativo:

\begin{sphinxVerbatim}[commandchars=\\\{\}]
\PYG{g+gp}{\PYGZdl{}} php prueba.php
\end{sphinxVerbatim}

\item[{Interactiva:}] \leavevmode
El intérprete interactivo solicita al usuario que introduzca una sentencia,
normalmente por teclado. Una vez introducida, el intérprete la analiza, la
ejecuta y vuelve a solicitar al usuario la introducción de una nueva
sentencia. El usuario, por tanto, ve inmediatamente el efecto que produce la
ejecución de la sentencia que acaba de introducir en el intérprete
interactivo.

\end{description}

\textless{}\sphinxurl{http://php.net/manual/es/language.basic-syntax.phpmode.php}\textgreater{}


\subsection{Intérprete interactivo}
\label{\detokenize{php:interprete-interactivo}}
La meta principal de este curso es escribir programas y, por tanto, la ejecución
por lotes es la más importante y la razón de ser del lenguaje. Pero el
intérprete interactivo resulta muy útil para hacerse con el manejo básico del
lenguaje y para realizar pruebas rápidas sin necesidad de tener que escribir un
programa expresamente para tal fin. Por ello, su utilidad didáctica es
innegable, así que empezaremos con él.


\subsubsection{Intérprete interactivo integrado (\sphinxstyleliteralintitle{php -a})}
\label{\detokenize{php:interprete-interactivo-integrado-php-a}}
PHP dispone de un intérprete interactivo integrado que, si bien es algo
espartano y no dispone de características adicionales que sí podemos encontrar
en otros intérpretes más avanzados, tiene lo justo y necesario para cumplir su
función.

Para empezar a trabajar con dicho intérprete de forma interactiva, usaremos el
comando \sphinxcode{php} con la opción \sphinxcode{-a} desde la consola del sistema operativo:

\begin{sphinxVerbatim}[commandchars=\\\{\}]
\PYG{g+gp}{\PYGZdl{}} php \PYGZhy{}a
\PYG{g+go}{Interactive mode enabled}

\PYG{g+go}{php \PYGZgt{}}
\end{sphinxVerbatim}

El intérprete nos muestra el \sphinxstyleemphasis{prompt} \sphinxcode{php \textgreater{}}, indicándonos que está listo
para recibir nuestras sentencias PHP. Probamos:

\begin{sphinxVerbatim}[commandchars=\\\{\}]
\PYG{g+go}{php \PYGZgt{} echo 75;}
\PYG{g+go}{75}
\PYG{g+go}{php \PYGZgt{} echo 23 * 5;}
\PYG{g+go}{115}
\PYG{g+go}{php \PYGZgt{} echo \PYGZdq{}Hola a todos\PYGZdq{};}
\PYG{g+go}{Hola a todos}
\PYG{g+go}{php \PYGZgt{} var\PYGZus{}dump(3 + 5);}
\PYG{g+go}{int(8)}
\PYG{g+go}{php \PYGZgt{}}
\end{sphinxVerbatim}

Como se ve, el intérprete ejecuta inmediatamente el comando introducido,
llevando a cabo las operaciones indicadas en la instrucción (en este caso,
evaluar la expresión y mostrar el resultado en pantalla) y, a continuación,
solicita un nuevo comando al usuario.

La sentencia introducida debe ser sintácticamente correcta. Si, por ejemplo, nos
olvidamos de escribir el punto y coma (\sphinxcode{;}), no obtendremos el resultado
esperado:

\begin{sphinxVerbatim}[commandchars=\\\{\}]
\PYG{g+go}{php \PYGZgt{} echo 73}
\PYG{g+go}{php \PYGZgt{} echo 25;}
\PYG{g+go}{PHP Parse error: syntax error, unexpected \PYGZsq{}echo\PYGZsq{} (T\PYGZus{}ECHO), expecting \PYGZsq{},\PYGZsq{} or}
\PYG{g+go}{\PYGZsq{};\PYGZsq{} in php shell code on line 2}

\PYG{g+go}{Parse error: syntax error, unexpected \PYGZsq{}echo\PYGZsq{} (T\PYGZus{}ECHO), expecting \PYGZsq{},\PYGZsq{} or \PYGZsq{};\PYGZsq{}}
\PYG{g+go}{in php shell code on line 2}

\PYG{g+go}{php \PYGZgt{}}
\end{sphinxVerbatim}

El mensaje de error se debe a que el intérprete ha considerado los dos comandos
\sphinxcode{echo} como si fueran una única sentencia:

\begin{sphinxVerbatim}[commandchars=\\\{\}]
\PYG{k}{echo} \PYG{l+m+mi}{73} \PYG{k}{echo} \PYG{l+m+mi}{25}\PYG{p}{;}
\end{sphinxVerbatim}

puesto que sólo ha encontrado un \sphinxcode{;} al final del todo. Por tanto, se queja de
que se ha encontrado la palabra \sphinxcode{echo} detrás del \sphinxcode{73} cuando se esperaba
una \sphinxcode{,} o un \sphinxcode{;}.

Por otra parte, si introducimos como sentencia una expresión acabada en \sphinxcode{;},
no obtendremos ningún resultado en pantalla, ya que la expresión se evaluará sin
más pero no se hará nada con dicho valor:

\begin{sphinxVerbatim}[commandchars=\\\{\}]
\PYG{g+go}{php \PYGZgt{} 3 + 5;}
\PYG{g+go}{php \PYGZgt{} 6 * 9;}
\PYG{g+go}{php \PYGZgt{}}
\end{sphinxVerbatim}

Lo que demuestra su nula utilidad práctica.

Para salir del intérprete interactivo, pulsamos la combinación de teclas \sphinxcode{Control-D}:

\begin{sphinxVerbatim}[commandchars=\\\{\}]
\PYG{g+go}{php \PYGZgt{} \PYGZca{}D}
\PYG{g+gp}{\PYGZdl{}}
\end{sphinxVerbatim}


\subsubsection{PsySH}
\label{\detokenize{php:psysh}}
\sphinxhref{http://psysh.org/}{PsySH}%
\begin{footnote}[7]\sphinxAtStartFootnote
\sphinxnolinkurl{http://psysh.org/}
%
\end{footnote} es una interesantísima aplicación desarrollada por
Justin Hileman (y otros) que proporciona un intérprete interactivo para PHP
bastante más potente y cómodo que el intérprete interactivo integrado que trae
PHP de serie. Entre sus características, incluye:
\begin{itemize}
\item {} 
Un \sphinxstyleemphasis{depurador} integrado que facilita la introspección de los programas y la
localización de errores.

\item {} 
Autocompletado pulsando \sphinxcode{Tab}.

\item {} 
Uso adecuado de espacios de nombres.

\item {} 
Histórico de órdenes introducidas.

\item {} 
Visualización a todo color.

\item {} 
Admite sentencias y expresiones.

\end{itemize}

La verdad es que, existiendo una herramienta así, no tiene demasiado sentido
usar el intérprete interactivo integrado de PHP. Tú simplemente haz la prueba,
comprueba la diferencia y dime si tengo razón o no…

La instalación de PsySH es muy sencilla:

\begin{sphinxVerbatim}[commandchars=\\\{\}]
\PYG{g+gp}{\PYGZdl{}} wget https://git.io/psysh
\PYG{g+gp}{\PYGZdl{}} chmod +x psysh
\PYG{g+gp}{\PYGZdl{}} sudo mv \PYGZhy{}f psysh /usr/local/bin
\end{sphinxVerbatim}

Con esto tenemos la herramienta básica. Si además queremos disponer del manual de PHP en línea (cosa altamente recomendable), hacemos también lo siguiente:

\begin{sphinxVerbatim}[commandchars=\\\{\}]
\PYG{g+gp}{\PYGZdl{}} wget https://psysh.org/manual/es/php\PYGZus{}manual.sqlite
\PYG{g+gp}{\PYGZdl{}} mkdir \PYGZhy{}p \PYGZti{}/.local/share/psysh
\PYG{g+gp}{\PYGZdl{}} mv \PYGZhy{}f php\PYGZus{}manual.sqlite \PYGZti{}/.local/share/psysh
\end{sphinxVerbatim}


\subsection{Modo dual de operación}
\label{\detokenize{php:modo-dual-de-operacion}}
ricpelo’s note: Se llaman \sphinxstyleemphasis{modo HTML} y \sphinxstyleemphasis{modo PHP}.


\subsection{Etiquetas \sphinxstyleliteralintitle{\textless{}?php} y \sphinxstyleliteralintitle{?\textgreater{}}}
\label{\detokenize{php:etiquetas-php-y}}

\section{Variables}
\label{\detokenize{php:variables}}
\textless{}\sphinxurl{http://php.net/manual/es/language.variables.php}\textgreater{}


\subsection{Conceptos básicos}
\label{\detokenize{php:conceptos-basicos}}
\textless{}\sphinxurl{http://php.net/manual/es/language.variables.basics.php}\textgreater{}


\subsection{Destrucción de variables}
\label{\detokenize{php:destruccion-de-variables}}
\textless{}\sphinxurl{http://php.net/manual/es/function.unset.php}\textgreater{}


\subsection{Operadores de asignación por valor y por referencia}
\label{\detokenize{php:operadores-de-asignacion-por-valor-y-por-referencia}}
\textless{}\sphinxurl{http://php.net/manual/es/language.operators.assignment.php}\textgreater{}

ricpelo’s note: En \sphinxcode{\$b =\& \$a;}, \sphinxcode{\$b} \sphinxstylestrong{NO} está apuntando a \sphinxcode{\$a}
o viceversa. Ambos apuntan al mismo
lugar. \textless{}\sphinxurl{http://php.net/manual/es/language.references.whatdo.php}\textgreater{}


\subsection{Variables predefinidas}
\label{\detokenize{php:variables-predefinidas}}
\textless{}\sphinxurl{http://php.net/manual/es/reserved.variables.php}\textgreater{}

ricpelo’s note: \sphinxcode{\$\_ENV} no funciona en la instalación actual (ver
\sphinxcode{variables\_order} en \sphinxcode{php.ini}. Habría que usar \sphinxcode{get\_env()}.


\section{Tipos básicos de datos}
\label{\detokenize{php:tipos-basicos-de-datos}}
\textless{}\sphinxurl{http://php.net/manual/es/language.types.intro.php}\textgreater{}


\subsection{Lógicos (\sphinxstyleliteralintitle{bool})}
\label{\detokenize{php:logicos-bool}}
\textless{}\sphinxurl{http://php.net/manual/es/language.types.boolean.php}\textgreater{}

\begin{DUlineblock}{0em}
\item[] ricpelo’s note: Se escriben en minúscula: \sphinxcode{false} y
\sphinxcode{true}. \textless{}\sphinxurl{https://github.com/yiisoft/yii2/blob/master/docs/internals/core-code-style.md\#51-types}\textgreater{}
\item[] ricpelo’s note: \sphinxcode{boolean} es sinónimo de \sphinxcode{bool}, pero debería
usarse \sphinxcode{bool}.
\end{DUlineblock}


\subsubsection{Operadores lógicos}
\label{\detokenize{php:operadores-logicos}}
\textless{}\sphinxurl{http://php.net/manual/es/language.operators.logical.php}\textgreater{}

\begin{DUlineblock}{0em}
\item[] ricpelo’s note: \sphinxstyleemphasis{Cuidado}:
\item[] - \sphinxcode{false and (true \&\& print('hola'))} no imprime nada y devuelve
\sphinxcode{false}, por lo que \sphinxstylestrong{el código va en cortocircuito y se evalúa de
izquierda a derecha} incluso aunque el \sphinxcode{\&\&} y los paréntesis tengan
más prioridad que el \sphinxcode{and}.
\item[] - Otra forma de verlo es comprobar que
\sphinxcode{print('uno') and (1 + print('dos'))} escribe \sphinxcode{unodos} (y devuelve
\sphinxcode{true}), por lo que la evaluación de los operandos del \sphinxcode{and} se
hace de izquierda a derecha aunque el \sphinxcode{+} tenga más prioridad (y
encima vaya entre paréntesis).
\item[] - En el \sphinxhref{http://php.net/manual/es/language.operators.precedence.php}{manual de
PHP}%
\begin{footnote}[8]\sphinxAtStartFootnote
\sphinxnolinkurl{http://php.net/manual/es/language.operators.precedence.php}
%
\end{footnote} se
dice que: \sphinxstyleemphasis{«La precedencia y asociatividad de los operadores solamente
determinan cómo se agrupan las expresiones, no especifican un orden de
evaluación. PHP no especifica (en general) el orden en que se evalúa
una expresión y se debería evitar el código que se asume un orden
específico de evaluación, ya que el comportamiento puede cambiar entre
versiones de PHP o dependiendo de código circundante.»}
\item[] - \sphinxhref{https://stackoverflow.com/questions/46861563/false-and-true-printhi}{Pregunta que hice al respecto en
StackOverflow}%
\begin{footnote}[9]\sphinxAtStartFootnote
\sphinxnolinkurl{https://stackoverflow.com/questions/46861563/false-and-true-printhi}
%
\end{footnote}.
\end{DUlineblock}


\subsection{Numéricos}
\label{\detokenize{php:numericos}}

\subsubsection{Enteros (\sphinxstyleliteralintitle{int})}
\label{\detokenize{php:enteros-int}}
\textless{}\sphinxurl{http://php.net/manual/es/language.types.integer.php}\textgreater{}

ricpelo’s note: \sphinxcode{integer} es sinónimo de \sphinxcode{int}, pero debería usarse
\sphinxcode{int}.


\subsubsection{Números en coma flotante (\sphinxstyleliteralintitle{float})}
\label{\detokenize{php:numeros-en-coma-flotante-float}}
\textless{}\sphinxurl{http://php.net/manual/es/language.types.float.php}\textgreater{}

ricpelo’s note: \sphinxcode{double} es sinónimo de \sphinxcode{float}, pero debería usarse
\sphinxcode{float}.


\subsubsection{Operadores}
\label{\detokenize{php:id4}}

\paragraph{Operadores aritméticos}
\label{\detokenize{php:operadores-aritmeticos}}
\textless{}\sphinxurl{http://php.net/manual/es/language.operators.arithmetic.php}\textgreater{}


\paragraph{Operadores de incremento/decremento}
\label{\detokenize{php:operadores-de-incremento-decremento}}
\textless{}\sphinxurl{http://php.net/manual/es/language.operators.increment.php}\textgreater{}


\subsection{Cadenas (\sphinxstyleliteralintitle{string})}
\label{\detokenize{php:cadenas-string}}
\textless{}\sphinxurl{http://php.net/manual/es/language.types.string.php}\textgreater{}

ricpelo’s note: Se usa \sphinxcode{\{\$var\}} y no
\sphinxcode{\$\{var\}} \textless{}\sphinxurl{https://github.com/yiisoft/yii2/blob/master/docs/internals/core-code-style.md\#variable-substitution}\textgreater{}


\subsubsection{Operadores de cadenas}
\label{\detokenize{php:operadores-de-cadenas}}
\textless{}\sphinxurl{http://php.net/manual/es/language.operators.string.php}\textgreater{}


\paragraph{Concatenación}
\label{\detokenize{php:concatenacion}}

\paragraph{Acceso y modificación por caracteres}
\label{\detokenize{php:acceso-y-modificacion-por-caracteres}}
\textless{}\sphinxurl{http://php.net/manual/es/language.types.string.php\#language.types.string.substr}\textgreater{}

\begin{DUlineblock}{0em}
\item[] ricpelo’s note: - \sphinxcode{echo \$a{[}3{]}}
\item[] - \sphinxcode{\$a{[}3{]} = 'x';}
\end{DUlineblock}


\paragraph{Operadores de incremento/decremento}
\label{\detokenize{php:id5}}
\textless{}\sphinxurl{http://php.net/manual/es/language.operators.increment.php}\textgreater{}


\subsubsection{Funciones de manejo de cadenas}
\label{\detokenize{php:funciones-de-manejo-de-cadenas}}
\textless{}\sphinxurl{http://php.net/ref.strings}\textgreater{}


\subsubsection{Extensión \sphinxstyleemphasis{mbstring}}
\label{\detokenize{php:extension-mbstring}}
\textless{}\sphinxurl{http://php.net/manual/en/book.mbstring.php}\textgreater{}

\begin{DUlineblock}{0em}
\item[] ricpelo’s note: - \sphinxcode{\$a{[}3{]}} equivale a \sphinxcode{mb\_substr(\$a, 3, 1)}
\item[] - \sphinxcode{\$a{[}3{]} = 'x';} no tiene equivalencia directa. Se podría hacer:
\item[] \sphinxcode{\$a = mb\_substr(\$a, 2, 1) . 'x' . mb\_substr(\$a, 4);}
\end{DUlineblock}


\subsection{Nulo}
\label{\detokenize{php:nulo}}
\textless{}\sphinxurl{http://php.net/manual/es/language.types.null.php}\textgreater{}

\begin{DUlineblock}{0em}
\item[] ricpelo’s note: \sphinxcode{{}`is\_null()} vs.
\sphinxcode{=== null} \textless{}\sphinxurl{https://phpbestpractices.org/\#checking-for-null}\textgreater{}
\item[] ricpelo’s note: El tipo \sphinxcode{null} y el valor \sphinxcode{null} se escriben en
minúscula. \textless{}\sphinxurl{https://github.com/yiisoft/yii2/blob/master/docs/internals/core-code-style.md\#51-types}\textgreater{}
\end{DUlineblock}


\subsection{Precedencia de operadores}
\label{\detokenize{php:precedencia-de-operadores}}
\textless{}\sphinxurl{http://php.net/manual/es/language.operators.precedence.php}\textgreater{}


\subsection{Operadores de asignación compuesta}
\label{\detokenize{php:operadores-de-asignacion-compuesta}}
ricpelo’s note: \sphinxcode{\$x} \sphinxstyleemphasis{\textless{}op\textgreater{}}\sphinxcode{= \$y}


\subsection{Comprobaciones}
\label{\detokenize{php:comprobaciones}}

\subsubsection{De tipos}
\label{\detokenize{php:de-tipos}}

\paragraph{\sphinxstyleliteralintitle{gettype()}}
\label{\detokenize{php:gettype}}
\textless{}\sphinxurl{http://php.net/manual/en/function.gettype.php}\textgreater{}


\paragraph{\sphinxstyleliteralintitle{is\_*()}}
\label{\detokenize{php:is}}
\textless{}\sphinxurl{http://php.net/manual/es/ref.var.php}\textgreater{}

ricpelo’s note: Poco útiles en formularios, ya que sólo se reciben
\sphinxcode{string}s.


\subsubsection{De valores}
\label{\detokenize{php:de-valores}}

\paragraph{\sphinxstyleliteralintitle{is\_numeric()}}
\label{\detokenize{php:is-numeric}}
\textless{}\sphinxurl{http://php.net/manual/es/function.is-numeric.php}\textgreater{}


\paragraph{\sphinxstyleliteralintitle{ctype\_*()}}
\label{\detokenize{php:ctype}}
\textless{}\sphinxurl{http://php.net/manual/es/book.ctype.php}\textgreater{}


\subsection{Conversiones}
\label{\detokenize{php:conversiones}}
\textless{}\sphinxurl{http://php.net/manual/es/language.types.type-juggling.php}\textgreater{}


\subsubsection{Coerción, moldeado, forzado o \sphinxstyleemphasis{casting}}
\label{\detokenize{php:coercion-moldeado-forzado-o-casting}}
\textless{}\sphinxurl{http://php.net/manual/es/language.types.type-juggling.php\#language.types.typecasting}\textgreater{}

ricpelo’s note: Conversión de cadena a número


\paragraph{Conversión a \sphinxstyleliteralintitle{bool}}
\label{\detokenize{php:conversion-a-bool}}
\textless{}\sphinxurl{http://php.net/manual/es/language.types.boolean.php\#language.types.boolean.casting}\textgreater{}


\paragraph{Conversión a \sphinxstyleliteralintitle{int}}
\label{\detokenize{php:conversion-a-int}}
\textless{}\sphinxurl{http://php.net/manual/es/language.types.integer.php\#language.types.integer.casting}\textgreater{}


\paragraph{Conversión a \sphinxstyleliteralintitle{float}}
\label{\detokenize{php:conversion-a-float}}
\textless{}\sphinxurl{http://php.net/manual/es/language.types.float.php\#language.types.float.casting}\textgreater{}


\paragraph{Conversión de \sphinxstyleliteralintitle{string} a número}
\label{\detokenize{php:conversion-de-string-a-numero}}
\textless{}\sphinxurl{http://php.net/manual/es/language.types.string.php\#language.types.string.conversion}\textgreater{}

ricpelo’s note: \sphinxstylestrong{¡Cuidado!}: La documentación dice que \sphinxcode{1 + "pepe"}
o \sphinxcode{1 + "10 pepe"} funciona, pero en PHP7.1 da un \sphinxstylestrong{PHP Warning: A
non-numeric value encountered}.


\paragraph{Conversión a \sphinxstyleliteralintitle{string}}
\label{\detokenize{php:conversion-a-string}}
\textless{}\sphinxurl{http://php.net/manual/es/language.types.string.php\#language.types.string.casting}\textgreater{}


\subsubsection{Funciones de obtención de valores}
\label{\detokenize{php:funciones-de-obtencion-de-valores}}
ricpelo’s note: Hacen más o menos lo mismo que los \sphinxstyleemphasis{casting} pero con
funciones en lugar de con operadores. Puede ser interesante porque las
funciones se pueden guardar, usar con \sphinxstyleemphasis{map}, \sphinxstyleemphasis{reduce}, etc.


\paragraph{\sphinxstyleliteralintitle{intval()}}
\label{\detokenize{php:intval}}
\textless{}\sphinxurl{http://php.net/manual/es/function.intval.php}\textgreater{}


\paragraph{\sphinxstyleliteralintitle{floatval()}}
\label{\detokenize{php:floatval}}
\textless{}\sphinxurl{http://php.net/manual/es/function.floatval.php}\textgreater{}


\paragraph{\sphinxstyleliteralintitle{strval()}}
\label{\detokenize{php:strval}}
\textless{}\sphinxurl{http://php.net/manual/es/function.strval.php}\textgreater{}


\paragraph{\sphinxstyleliteralintitle{boolval()}}
\label{\detokenize{php:boolval}}
\textless{}\sphinxurl{http://php.net/manual/es/function.boolval.php}\textgreater{}


\subsubsection{Funciones de formateado numérico}
\label{\detokenize{php:funciones-de-formateado-numerico}}

\paragraph{\sphinxstyleliteralintitle{number\_format()}}
\label{\detokenize{php:number-format}}
\textless{}\sphinxurl{http://php.net/manual/es/function.number-format.php}\textgreater{}


\paragraph{\sphinxstyleliteralintitle{money\_format()}}
\label{\detokenize{php:money-format}}
\textless{}\sphinxurl{http://php.net/manual/es/function.money-format.php}\textgreater{}


\subparagraph{\sphinxstyleliteralintitle{setlocale()}}
\label{\detokenize{php:setlocale}}
\textless{}\sphinxurl{http://php.net/manual/es/function.setlocale.php}\textgreater{}

ricpelo’s note:
\sphinxcode{setlocale(LC\_ALL, 'es\_ES.UTF-8'); // Hay que poner el *locale* completo, con la codificación y todo (.UTF-8)}


\subsection{Comparaciones}
\label{\detokenize{php:comparaciones}}

\subsubsection{Operadores de comparación}
\label{\detokenize{php:operadores-de-comparacion}}
\textless{}\sphinxurl{http://php.net/manual/es/language.operators.comparison.php}\textgreater{}


\subsubsection{\sphinxstyleliteralintitle{==} vs. \sphinxstyleliteralintitle{===}}
\label{\detokenize{php:vs}}

\subsubsection{Ternario (\sphinxstyleliteralintitle{?:})}
\label{\detokenize{php:ternario}}
\textless{}\sphinxurl{http://php.net/manual/es/language.operators.comparison.php\#language.operators.comparison.ternary}\textgreater{}


\subsubsection{Fusión de null (\sphinxstyleliteralintitle{??})}
\label{\detokenize{php:fusion-de-null}}
\textless{}\sphinxurl{https://wiki.php.net/rfc/isset\_ternary}\textgreater{}

ricpelo’s note: Equivalente al \sphinxcode{COALESCE()} de SQL.


\subsubsection{Reglas de comparación de tipos}
\label{\detokenize{php:reglas-de-comparacion-de-tipos}}
\textless{}\sphinxurl{http://php.net/manual/es/types.comparisons.php}\textgreater{}

ricpelo’s note: \sphinxcode{"250" \textless{} "27"} devuelve \sphinxcode{false}


\section{Constantes}
\label{\detokenize{php:constantes}}
\textless{}\sphinxurl{http://php.net/manual/es/language.constants.syntax.php}\textgreater{}

\begin{DUlineblock}{0em}
\item[] ricpelo’s note: Diferencias entre constantes y variables:
\item[] - Las constantes no llevan el signo dólar (\sphinxcode{\$}) como prefijo.
\item[] - Antes de PHP 5.3, las constantes solo podían ser definidas usando la
función \sphinxcode{define()} y no por simple asignación.
\item[] - Las constantes pueden ser definidas y accedidas desde cualquier
sitio sin importar las reglas de acceso de variables.
\item[] - Las constantes no pueden ser redefinidas o eliminadas una vez se han
definido.
\item[] - Las constantes podrían evaluarse como valores escalares. A partir de
PHP 5.6 es posible definir una constante de array con la palabra
reservada \sphinxcode{const}, y, a partir de PHP 7, las constantes de array
también se pueden definir con \sphinxcode{define()}. Se pueden utilizar arrays
en expresiones escalares constantes (por ejemplo,
\sphinxcode{const FOO = array(1,2,3){[}0{]};}), aunque el resultado final debe ser
un valor de un tipo permitido.
\end{DUlineblock}


\subsection{\sphinxstyleliteralintitle{define()} y \sphinxstyleliteralintitle{const}}
\label{\detokenize{php:define-y-const}}

\subsection{Constantes predefinidas}
\label{\detokenize{php:constantes-predefinidas}}
\textless{}\sphinxurl{http://php.net/manual/es/language.constants.predefined.php}\textgreater{}


\subsection{\sphinxstyleliteralintitle{defined()}}
\label{\detokenize{php:defined}}
\textless{}\sphinxurl{http://php.net/manual/es/function.defined.php}\textgreater{}


\section{Flujo de control}
\label{\detokenize{php:flujo-de-control}}

\subsection{Estructuras de control}
\label{\detokenize{php:estructuras-de-control}}
\textless{}\sphinxurl{http://php.net/manual/es/language.control-structures.php}\textgreater{}


\subsubsection{Sintaxis alternativa}
\label{\detokenize{php:sintaxis-alternativa}}
\textless{}\sphinxurl{http://php.net/manual/es/control-structures.alternative-syntax.php}\textgreater{}

ricpelo’s note: El \sphinxcode{do \{ ... \} while (...);} \sphinxstylestrong{no} tiene sintaxis
alternativa.


\subsection{Inclusión de archivos}
\label{\detokenize{php:inclusion-de-archivos}}

\subsubsection{\sphinxstyleliteralintitle{include}, \sphinxstyleliteralintitle{require}}
\label{\detokenize{php:include-require}}
\textless{}\sphinxurl{http://php.net/manual/es/function.include.php}\textgreater{}

\begin{DUlineblock}{0em}
\item[] ricpelo’s note: El nombre del archivo debe aparecer con su extensión.
No vale hacer \sphinxcode{require 'pepe';}.
\item[] ricpelo’s note: Cuando un archivo es incluido, el intérprete abandona
el modo PHP e ingresa al modo HTML al comienzo del archivo objetivo y
se reanuda de nuevo al final.
\item[] ricpelo’s note: Si el archivo incluido tiene un \sphinxcode{return ...;}, el
\sphinxcode{include} o el \sphinxcode{require} que lo incluya devolverá el valor
devuelto por el \sphinxcode{return}.
\end{DUlineblock}


\subsubsection{\sphinxstyleliteralintitle{include\_once}, \sphinxstyleliteralintitle{require\_once}}
\label{\detokenize{php:include-once-require-once}}
\textless{}\sphinxurl{http://php.net/manual/es/function.include-once.php}\textgreater{}


\section{Funciones predefinidas destacadas}
\label{\detokenize{php:funciones-predefinidas-destacadas}}

\subsection{\sphinxstyleliteralintitle{isset()}}
\label{\detokenize{php:isset}}
\textless{}\sphinxurl{http://php.net/manual/es/function.isset.php}\textgreater{}

\begin{DUlineblock}{0em}
\item[] ricpelo’s note: Cuidado si la variable contiene \sphinxcode{null}.
\item[] ricpelo’s note: No da error ni advertencia si la variable no existe.
\end{DUlineblock}


\subsection{\sphinxstyleliteralintitle{empty()}}
\label{\detokenize{php:empty}}
\textless{}\sphinxurl{http://php.net/manual/es/function.empty.php}\textgreater{}

ricpelo’s note: Para evitar el problema de \sphinxcode{empty("0") === true}:

\begin{sphinxVerbatim}[commandchars=\\\{\}]
\PYG{k}{function} \PYG{n+nf}{is\PYGZus{}blank}\PYG{p}{(}\PYG{n+nv}{\PYGZdl{}value}\PYG{p}{)} \PYG{p}{\PYGZob{}}
    \PYG{k}{return} \PYG{k}{empty}\PYG{p}{(}\PYG{n+nv}{\PYGZdl{}value}\PYG{p}{)} \PYG{o}{\PYGZam{}\PYGZam{}} \PYG{o}{!}\PYG{n+nb}{is\PYGZus{}numeric}\PYG{p}{(}\PYG{n+nv}{\PYGZdl{}value}\PYG{p}{);}
\PYG{p}{\PYGZcb{}}
\end{sphinxVerbatim}

ricpelo’s note: No da error ni advertencia si la variable no existe.


\subsection{\sphinxstyleliteralintitle{var\_dump()}}
\label{\detokenize{php:id6}}
\textless{}\sphinxurl{http://php.net/manual/es/function.var-dump.php}\textgreater{}


\section{Arrays}
\label{\detokenize{php:arrays}}
\textless{}\sphinxurl{http://php.net/manual/es/language.types.array.php}\textgreater{}

ricpelo’s note: Las claves pueden ser enteros o cadenas.


\subsection{Operadores para arrays}
\label{\detokenize{php:operadores-para-arrays}}
\textless{}\sphinxurl{http://php.net/manual/es/language.operators.array.php}\textgreater{}

ricpelo’s note: \sphinxstylestrong{Comparaciones}: Un \sphinxcode{array} con menos elementos es
menor. De otra forma, compara valor por valor.


\subsubsection{Acceso, modificación y agregación}
\label{\detokenize{php:acceso-modificacion-y-agregacion}}
\textless{}\sphinxurl{http://php.net/manual/es/language.types.array.php\#language.types.array.syntax.modifying}\textgreater{}


\subsection{Funciones de manejo de arrays{]}}
\label{\detokenize{php:funciones-de-manejo-de-arrays}}
\textless{}\sphinxurl{http://php.net/manual/es/book.array.php}\textgreater{}
\textless{}\sphinxurl{http://php.net/manual/es/ref.array.php}\textgreater{}


\subsubsection{Ordenación de arrays}
\label{\detokenize{php:ordenacion-de-arrays}}
\textless{}\sphinxurl{http://php.net/manual/es/array.sorting.php}\textgreater{}


\subsubsection{\sphinxstyleliteralintitle{print\_r()}}
\label{\detokenize{php:print-r}}

\subsubsection{\sphinxstyleliteralintitle{'+'} vs. \sphinxstyleliteralintitle{array\_merge()}}
\label{\detokenize{php:vs-array-merge}}

\subsubsection{\sphinxstyleliteralintitle{isset()} vs. \sphinxstyleliteralintitle{array\_key\_exists()}}
\label{\detokenize{php:isset-vs-array-key-exists}}
\textless{}\sphinxurl{http://php.net/manual/es/function.array-key-exists.php\#107786}\textgreater{}


\subsection{\sphinxstyleliteralintitle{foreach}}
\label{\detokenize{php:foreach}}
\textless{}\sphinxurl{http://php.net/manual/es/control-structures.foreach.php}\textgreater{}


\subsection{Conversión a \sphinxstyleliteralintitle{array}}
\label{\detokenize{php:conversion-a-array}}
\textless{}\sphinxurl{http://php.net/manual/es/language.types.array.php\#language.types.array.casting}\textgreater{}


\subsection{\sphinxstyleemphasis{Ejemplo}: \sphinxstyleliteralintitle{\$argv} en CLI}
\label{\detokenize{php:ejemplo-argv-en-cli}}
\textless{}\sphinxurl{http://php.net/manual/es/reserved.variables.argv.php}\textgreater{}


\section{Funciones definidas por el usuario}
\label{\detokenize{php:funciones-definidas-por-el-usuario}}
\textless{}\sphinxurl{http://php.net/manual/es/language.functions.php}\textgreater{}


\subsection{Argumentos}
\label{\detokenize{php:argumentos}}
\textless{}\sphinxurl{http://php.net/manual/es/functions.arguments.php}\textgreater{}


\subsubsection{Paso de argumentos por valor y por referencia}
\label{\detokenize{php:paso-de-argumentos-por-valor-y-por-referencia}}
\textless{}\sphinxurl{http://php.net/manual/es/functions.arguments.php\#functions.arguments.by-reference}\textgreater{}


\subsubsection{Argumentos por defecto}
\label{\detokenize{php:argumentos-por-defecto}}
\textless{}\sphinxurl{http://php.net/manual/es/functions.arguments.php\#functions.arguments.default}\textgreater{}

ricpelo’s note:
\sphinxcode{php   function prueba(\$opciones = {[}{]}) \{       extract(\$opciones);       // ...   \}}


\subsection{Ámbito de variables}
\label{\detokenize{php:ambito-de-variables}}
\textless{}\sphinxurl{http://php.net/language.variables.scope}\textgreater{}


\subsubsection{Ámbito simple al archivo}
\label{\detokenize{php:ambito-simple-al-archivo}}

\subsubsection{Variables locales}
\label{\detokenize{php:variables-locales}}

\subsubsection{Uso de \sphinxstyleliteralintitle{global}}
\label{\detokenize{php:uso-de-global}}
ricpelo’s note: Usar \sphinxcode{global \$x;} cuando \sphinxcode{\$x} no existe hace que
\sphinxcode{\$x} empiece a existir y valga \sphinxcode{null}.


\subsubsection{Variables superglobales}
\label{\detokenize{php:variables-superglobales}}
\textless{}\sphinxurl{http://php.net/manual/es/language.variables.superglobals.php}\textgreater{}


\subsection{Declaraciones de tipos}
\label{\detokenize{php:declaraciones-de-tipos}}
ricpelo’s note: \sphinxstylestrong{NO} se hacen conversiones implícitas a \sphinxcode{array}, ni
en argumentos ni en devolución.


\subsubsection{Declaraciones de tipo de argumento}
\label{\detokenize{php:declaraciones-de-tipo-de-argumento}}
\textless{}\sphinxurl{http://php.net/manual/es/functions.arguments.php\#functions.arguments.type-declaration}\textgreater{}


\subsubsection{Declaraciones de tipo de devolución}
\label{\detokenize{php:declaraciones-de-tipo-de-devolucion}}
\textless{}\sphinxurl{http://php.net/manual/es/functions.returning-values.php\#functions.returning-values.type-declaration}\textgreater{}


\subsubsection{Tipos \sphinxstyleemphasis{nullable} (\sphinxstyleliteralintitle{?}) y \sphinxstyleliteralintitle{void}}
\label{\detokenize{php:tipos-nullable-y-void}}
\textless{}\sphinxurl{http://php.net/manual/es/migration71.new-features.php}\textgreater{}


\subsubsection{Tipificación estricta}
\label{\detokenize{php:tipificacion-estricta}}
\textless{}\sphinxurl{http://php.net/manual/es/functions.arguments.php\#functions.arguments.type-declaration.strict}\textgreater{}

ricpelo’s note: El \sphinxcode{declare(strict\_types=1);} se pone en el archivo
que hace la llamada, no en el que define la función.


\part{Yii2}
\label{\detokenize{yii2:yii2}}\label{\detokenize{yii2::doc}}

\chapter{Introducción a Yii2}
\label{\detokenize{yii2:introduccion-a-yii2}}

\chapter{Conceptos fundamentales de Yii2}
\label{\detokenize{yii2:conceptos-fundamentales-de-yii2}}

\chapter{Estructura de una aplicación Yii2}
\label{\detokenize{yii2:estructura-de-una-aplicacion-yii2}}

\chapter{Gestión de peticiones en Yii2}
\label{\detokenize{yii2:gestion-de-peticiones-en-yii2}}

\chapter{Acceso a bases de datos en Yii2}
\label{\detokenize{yii2:acceso-a-bases-de-datos-en-yii2}}

\chapter{Recogida de datos y validación de formularios}
\label{\detokenize{yii2:recogida-de-datos-y-validacion-de-formularios}}

\chapter{Visualización de datos en Yii2}
\label{\detokenize{yii2:visualizacion-de-datos-en-yii2}}

\chapter{Características adicionales de Yii2}
\label{\detokenize{yii2:caracteristicas-adicionales-de-yii2}}

\chapter{Seguridad y cacheado en Yii2}
\label{\detokenize{yii2:seguridad-y-cacheado-en-yii2}}

\chapter{Pruebas, documentación y mantenimiento}
\label{\detokenize{yii2:pruebas-documentacion-y-mantenimiento}}

\chapter{Computación en la nube}
\label{\detokenize{yii2:computacion-en-la-nube}}

\chapter{Contenedores}
\label{\detokenize{yii2:contenedores}}

\part{Índices y tablas}
\label{\detokenize{index_latex:indices-y-tablas}}


\renewcommand{\indexname}{Índice}
\printindex
\end{document}